\documentclass[12pt,a4paper]{article}

% ===== PACKAGES =====
\usepackage[utf8]{inputenc}
\usepackage[vietnamese]{babel}
\usepackage{amsmath, amssymb, amsfonts}
\usepackage{geometry}
\usepackage{enumitem}
\usepackage{array}
\usepackage{booktabs}
\usepackage{fancyhdr}
\usepackage{titlesec}
\usepackage{xcolor}
\usepackage{tcolorbox}

\geometry{a4paper, margin=2.5cm}

% ===== CUSTOM STYLES (giống file mẫu đã verify) =====
\definecolor{mainblue}{RGB}{0,90,160}
\definecolor{darkgreen}{RGB}{0,120,60}
\definecolor{resultbg}{RGB}{230,245,255}

\newtcolorbox{resultbox}{
  colback=resultbg, colframe=mainblue,
  boxrule=0.8pt, arc=3pt,
  left=8pt, right=8pt, top=6pt, bottom=6pt,
  fonttitle=\bfseries, title=Kết quả
}

\newtcolorbox{notebox}[1][]{
  colback=yellow!8, colframe=orange!70!black,
  boxrule=0.6pt, arc=3pt,
  left=8pt, right=8pt, top=6pt, bottom=6pt,
  fonttitle=\bfseries, title={Lưu ý}, #1
}

\pagestyle{fancy}
\fancyhf{}
\fancyhead[L]{\textit{Phân tích \& Thiết kế Thuật toán}}
\fancyhead[R]{\textit{GV: Huỳnh Thị Thanh Thương}}
\fancyfoot[C]{\thepage}

\title{\textbf{Hướng dẫn giải HW\#02: Độ phức tạp \& Ký hiệu tiệm cận}\\[6pt]
\large (Có lập luận/giải thích theo hướng dẫn giảng viên)}
\author{}
\date{}

\begin{document}
\maketitle
\thispagestyle{fancy}

% ===== PHƯƠNG PHÁP CHUNG =====
\section*{Phương pháp giải (theo hướng dẫn của giảng viên)}

Theo cô Thương hướng dẫn trong buổi 4 và buổi 5:

\begin{enumerate}[label=\textbf{(\arabic*)}]
  \item \textbf{Không đoán mò} --- phải có biến đổi toán học làm cơ sở.
  \item \textbf{Đưa về cùng cơ số} --- nếu hai hàm đều ở dạng mũ nhưng khác cơ số, chuyển hết về cơ số nhỏ nhất (thường là $2$) rồi so phần mũ.
  \item \textbf{Phân lớp (phân nhóm)} --- quy mỗi hàm về lớp độ phức tạp quen thuộc ($\Theta$ tương ứng), rồi chỉ cần so sánh giữa các lớp.
  \item \textbf{Hệ số bỏ đi} --- trong cùng lớp, hệ số hằng là không đáng kể; chỉ giữ lại phần ``bậc'' cao nhất.
  \item \textbf{Kết luận}: ghi rõ hàm thuộc lớp $O/\Theta$ nào, không cần ghi đúng công thức gốc.
\end{enumerate}

\begin{notebox}
Thứ tự các lớp cơ bản cần nhớ:
\[
  \underbrace{\Theta(1)}_{\text{hằng}} \prec
  \underbrace{\Theta(\log n)}_{\text{logarit}} \prec
  \underbrace{\Theta(n^c)}_{\text{đa thức}} \prec
  \underbrace{\Theta(c^n)}_{\text{hàm mũ}} \prec
  \underbrace{\Theta(n!)}_{\text{giai thừa}} \prec
  \underbrace{\Theta(c^{c^n})}_{\text{siêu mũ}}
\]
Trong cùng lớp đa thức: so bậc. Trong cùng lớp mũ: đưa về cùng cơ số, so phần mũ.
\end{notebox}

\bigskip

\textbf{Tài liệu tham khảo:}
\begin{itemize}[noitemsep]
  \item {[1]} Anany Levitin, \textit{Introduction to the Design and Analysis of Algorithms}, 3rd Ed., 2014
  \item {[2]} Thomas H. Cormen, \textit{Introduction to Algorithms} (CLRS), 3rd Ed., MIT Press, 2009
\end{itemize}

% =========================================
%                BÀI 1
% =========================================
\newpage
\section{Bài 1: Sắp xếp hàm theo thứ tự tăng trưởng}

\subsection{Nhóm 1 (Group 1 --- 7 hàm)}

\textbf{Đề bài:} Sắp xếp các hàm sau theo bậc tăng trưởng từ thấp đến cao:
\[
  (n-2)!, \quad 5\lg(n+100)^{10}, \quad 2^{2n}, \quad 0.001n^4 + 3n^3 + 1, \quad \ln^2 n, \quad \sqrt[3]{n}, \quad 3^n.
\]

\subsubsection{Bước 1: Rút gọn và phân lớp từng hàm}

\begin{enumerate}[label=\textbf{\arabic*.}]

  \item $5\lg(n+100)^{10}$

  Theo tính chất logarit: $\lg(n+100)^{10} = 10\lg(n+100)$.

  Vậy: $5\lg(n+100)^{10} = 50\lg(n+100)$.

  Khi $n$ đủ lớn: $\lg(n+100) \approx \lg n$ (sai biệt không đáng kể).

  $\Longrightarrow$ Hàm này thuộc lớp $\Theta(\log n)$.

  \item $\ln^2 n$

  Tức là $(\ln n)^2$. Đây là hàm logarit bình phương.

  $\Longrightarrow$ Thuộc lớp $\Theta(\log^2 n)$.

  \item $\sqrt[3]{n}$

  Viết lại dưới dạng số mũ: $\sqrt[3]{n} = n^{1/3}$.

  $\Longrightarrow$ Thuộc lớp $\Theta(n^{1/3})$ --- đây là lớp \textbf{đa thức}, bậc $1/3$.

  \item $0.001n^4 + 3n^3 + 1$

  Là đa thức bậc 4 theo $n$. Bỏ hệ số, bỏ các bậc thấp hơn, chỉ giữ bậc cao nhất.

  $\Longrightarrow$ Thuộc lớp $\Theta(n^4)$.

  \item $3^n$

  \textbf{Chuyển về cơ số 2} (theo phương pháp cô hướng dẫn):
  \[
    3^n = \bigl(2^{\lg 3}\bigr)^n = 2^{n\lg 3} \approx 2^{1{,}585\,n}.
  \]
  Phần mũ xấp xỉ $1{,}585\,n$. $\Longrightarrow$ Thuộc lớp hàm mũ $\Theta(3^n)$.

  \item $2^{2n}$

  Rút gọn: $2^{2n} = (2^2)^n = 4^n = 2^{2n}$.

  \textbf{Chuyển về cơ số 2}: phần mũ là $2n$.

  So với $3^n$ (phần mũ $\approx 1{,}585\,n$): ta có $2n > 1{,}585\,n$ nên $4^n > 3^n$.

  $\Longrightarrow$ Thuộc lớp hàm mũ $\Theta(4^n)$.

  \item $(n-2)!$

  Đây là hàm giai thừa. Khi $n$ đủ lớn, $(n-2)!$ tăng nhanh hơn mọi hàm mũ.

  $\Longrightarrow$ Thuộc lớp $\Theta(n!)$.

\end{enumerate}

\subsubsection{Bước 2: Lập luận so sánh giữa các lớp}

\begin{enumerate}[label=\textbf{(\alph*)}]
  \item \textbf{Logarit $\prec$ Logarit bình phương:}
    $\log n \prec \log^2 n$ (bậc 1 < bậc 2 của log).

    $\Rightarrow 5\lg(n+100)^{10} \prec \ln^2 n$.

  \item \textbf{Logarit $\prec$ Đa thức:}
    Mọi hàm logarit (dù lũy thừa bao nhiêu) luôn tăng chậm hơn bất kỳ hàm đa thức nào (dù bậc nhỏ tới đâu).

    $\Rightarrow \ln^2 n \prec \sqrt[3]{n}$.

  \item \textbf{Đa thức bậc thấp $\prec$ Đa thức bậc cao:}
    So bậc: $1/3 < 4$.

    $\Rightarrow \sqrt[3]{n} \prec 0.001n^4 + 3n^3 + 1$.

  \item \textbf{Đa thức $\prec$ Hàm mũ:}
    Hàm đa thức (dù bậc cao tới đâu) luôn tăng chậm hơn hàm mũ.

    $\Rightarrow n^4 \prec 3^n$.

  \item \textbf{Hàm mũ cơ số nhỏ $\prec$ Hàm mũ cơ số lớn:}
    Đưa về cùng cơ số 2: phần mũ $1{,}585\,n < 2n$ nên $3^n \prec 4^n$.

    $\Rightarrow 3^n \prec 2^{2n}$.

  \item \textbf{Hàm mũ $\prec$ Giai thừa:}
    Hàm giai thừa tăng nhanh hơn bất kỳ hàm mũ nào.

    $\Rightarrow 2^{2n} \prec (n-2)!$.
\end{enumerate}

\begin{resultbox}
\textbf{Kết quả sắp xếp Nhóm 1} (từ thấp đến cao):
\[
  \boxed{
    5\lg(n+100)^{10}
    \;\prec\; \ln^2 n
    \;\prec\; \sqrt[3]{n}
    \;\prec\; 0.001n^4 + 3n^3 + 1
    \;\prec\; 3^n
    \;\prec\; 2^{2n}
    \;\prec\; (n-2)!
  }
\]
\end{resultbox}

\subsection{Nhóm 2 (Group 2 --- 30 hàm, bài 3-3 CLRS)}

\textbf{Đề bài:} Sắp xếp 30 hàm theo bậc tăng trưởng và phân thành các lớp tương đương
($f \equiv g$ khi $f = \Theta(g)$).

\textbf{Quy ước:} $\lg n = \log_2 n$, \; $\ln n = \log_e n$, \; $\lg^k n = (\lg n)^k$, \; $\lg\lg n = \lg(\lg n)$.

\subsubsection{Bước 1: Phân lớp từng hàm}

\textbf{Lớp 1: Hằng số --- $\Theta(1)$}

Các hàm: $1$ và $n^{1/\lg n}$.

Lập luận cho $n^{1/\lg n}$: Đặt $y = n^{1/\lg n}$. Lấy $\lg$ hai vế:
\[
  \lg y = \frac{1}{\lg n} \cdot \lg n = 1
  \quad\Longrightarrow\quad y = 2^1 = 2 \in \Theta(1).
\]

\textbf{Lớp 2: Logarit lặp --- $\Theta(\lg^* n)$ và lân cận}

Các hàm: $\lg(\lg^* n)$, $\lg^* n$, $\lg^*(\lg n)$, $2^{\lg^* n}$.

Lập luận:
\begin{itemize}[noitemsep]
  \item $\lg^*(\lg n) = \lg^* n - 1 \in \Theta(\lg^* n)$, nên $\lg^* n \equiv \lg^*(\lg n)$.
  \item $\lg(\lg^* n)$: lấy log của hàm log lặp $\Rightarrow$ tăng chậm hơn $\lg^* n$.
  \item $2^{\lg^* n}$: hàm mũ của $\lg^* n$ $\Rightarrow$ tăng nhanh hơn $\lg^* n$.
\end{itemize}

Thứ tự: $\lg(\lg^* n) \prec \lg^* n \equiv \lg^*(\lg n) \prec 2^{\lg^* n}$.

\textbf{Lớp 3: Double logarit --- $\Theta(\lg\lg n)$, $\Theta(\sqrt{\lg n})$}

$\lg\lg n \prec \sqrt{\lg n}$ vì $\lg\lg n$ tăng chậm hơn $(\lg n)^{1/2}$.

\textbf{Lớp 4: Logarit --- $\Theta(\ln n)$}

Các hàm: $\ln n$, $\lg n$. Khác cơ số chỉ là hệ số hằng: $\lg n = \frac{\ln n}{\ln 2}$.

$\Longrightarrow \ln n \equiv \lg n$.

\textbf{Lớp 5: Logarit bình phương --- $\Theta(\lg^2 n)$}

$(\lg n)^2$. Thuộc lớp logarit nhưng bậc cao hơn $\ln n$.

\textbf{Lớp 6: Căn --- $\Theta(\sqrt{n})$}

Các hàm: $(\sqrt{2})^{\lg n}$, $\sqrt{n}$.

Lập luận: $(\sqrt{2})^{\lg n} = 2^{(\lg n)/2} = (2^{\lg n})^{1/2} = n^{1/2} = \sqrt{n}$.

$\Longrightarrow (\sqrt{2})^{\lg n} \equiv \sqrt{n}$.

\textbf{Lớp 7: Tuyến tính --- $\Theta(n)$}

Các hàm: $n$, $2^{\lg n}$.

Lập luận: $2^{\lg n} = n$ (theo định nghĩa $\lg = \log_2$).

\textbf{Lớp 8: $n\log n$ --- $\Theta(n\lg n)$}

Các hàm: $n\lg n$, $\lg(n!)$.

Lập luận (Stirling): $\lg(n!) = \Theta(n\lg n)$.

\textbf{Lớp 9: Đa thức --- $\Theta(n^2)$, $\Theta(n^3)$, $\Theta(n^{\lg\lg n})$}

\begin{itemize}[noitemsep]
  \item $4^{\lg n} = 2^{2\lg n} = (2^{\lg n})^2 = n^2 \Longrightarrow 4^{\lg n} \equiv n^2$.
  \item $n^3$: đa thức bậc 3.
  \item $(\lg n)^{\lg n}$: Đặt $y = (\lg n)^{\lg n}$, lấy $\lg$: $\lg y = \lg n \cdot \lg\lg n$. Mà $n^{\lg\lg n}$: $\lg(n^{\lg\lg n}) = \lg\lg n \cdot \lg n$. Vậy $(\lg n)^{\lg n} \equiv n^{\lg\lg n}$.
\end{itemize}

Thứ tự: $n^2 \prec n^3 \prec n^{\lg\lg n}$ (vì $\lg\lg n \to \infty$).

\textbf{Lớp 10: Hàm mũ và siêu mũ}

Thứ tự (đưa về cơ số 2 so phần mũ):
\[
  (3/2)^n \prec 2^n \equiv 2^{n+1} \prec e^n \prec n! \prec (n+1)! \prec 2^{2^n} \prec 2^{2^{n+1}}
\]

Lập luận: $2^{n+1} = 2 \cdot 2^n \in \Theta(2^n)$. $(n+1)! = (n+1) \cdot n!$. $2^{2^{n+1}} = (2^{2^n})^2$.

\subsubsection{Bước 2: Tổng kết}

\begin{resultbox}
\textbf{Kết quả sắp xếp Nhóm 2} (từ thấp đến cao):
\begin{align*}
  1 &\prec \lg(\lg^* n) \prec \lg^* n \prec 2^{\lg^* n} \prec \lg\lg n \prec \sqrt{\lg n} \prec \ln n \\
    &\prec \lg^2 n \prec \sqrt{n} \prec n \prec n\lg n \prec n^2 \prec n^3 \prec n^{\lg\lg n} \\
    &\prec (3/2)^n \prec 2^n \prec e^n \prec n! \prec (n+1)! \prec 2^{2^n} \prec 2^{2^{n+1}}
\end{align*}
\end{resultbox}

% =========================================
%                BÀI 2
% =========================================
\newpage
\section{Bài 2: Xác định $O$, $\Theta$, $\Omega$ và Phân lớp hàm}

\subsection{Câu 1: Xác định đúng hay sai}

\textbf{Đề bài:} Dùng định nghĩa không hình thức (informal) của $O$, $\Theta$, $\Omega$ để xác định các khẳng định sau đúng hay sai.

\bigskip

Trước hết, ta rút gọn hàm cần xét:
\[
  \frac{n(n+1)}{2} = \frac{n^2 + n}{2} \in \Theta(n^2).
\]
Bỏ hệ số $1/2$, bỏ bậc thấp $n$, phần chính là $n^2$. Vậy bậc tăng trưởng là $n^2$.

\bigskip

\begin{enumerate}[label=\textbf{\alph*.}]

  \item $\dfrac{n(n+1)}{2} \in O(n^3)$ \hfill \textcolor{darkgreen}{\textbf{ĐÚNG}} \\[4pt]
  \textbf{Lập luận:} Ta đã biết $\frac{n(n+1)}{2} \in \Theta(n^2)$.
  Mà $n^2 \prec n^3$ (bậc $2 < 3$), nên $n^2 \in O(n^3)$.

  Theo tính chất bắc cầu: $\frac{n(n+1)}{2} \in O(n^2) \subseteq O(n^3)$.

  (Như cô giảng trong buổi 4: một hàm có thể thuộc nhiều lớp $O$ --- $O(n^2)$, $O(n^3)$, $O(n^4)$,\ldots\ đều đúng.
  Hàm bậc 2 bị chặn trên bởi $n^3$ là hiển nhiên.)

  \bigskip

  \item $\dfrac{n(n+1)}{2} \in O(n^2)$ \hfill \textcolor{darkgreen}{\textbf{ĐÚNG}} \\[4pt]
  \textbf{Lập luận:} $\frac{n(n+1)}{2} = \frac{n^2 + n}{2} \leq \frac{n^2 + n^2}{2} = n^2$ với mọi $n \geq 1$.

  Chọn $c = 1$, $n_0 = 1$: ta có $\frac{n(n+1)}{2} \leq 1 \cdot n^2$ với mọi $n \geq 1$.

  $\Longrightarrow$ Theo định nghĩa $O$: $\frac{n(n+1)}{2} \in O(n^2)$. \checkmark

  \bigskip

  \item $\dfrac{n(n+1)}{2} \in \Theta(n^3)$ \hfill \textcolor{red}{\textbf{SAI}} \\[4pt]
  \textbf{Lập luận (phản chứng --- theo cô hướng dẫn buổi 5):}

  Giả sử khẳng định đúng, tức $\frac{n(n+1)}{2} \in \Theta(n^3)$.

  Theo định nghĩa $\Theta$: $\exists\, c_1, c_2 > 0,\; n_0 \in \mathbb{N}$ sao cho:
  \[
    c_1 \cdot n^3 \leq \frac{n(n+1)}{2} \leq c_2 \cdot n^3 \quad \forall\, n \geq n_0.
  \]

  Xét vế trái. Bất đẳng thức $c_1 \cdot n^3 \leq \frac{n^2+n}{2}$ suy ra:
  \[
    c_1 \leq \frac{n^2+n}{2n^3} = \frac{1}{2n} + \frac{1}{2n^2} \to 0 \quad \text{khi } n \to \infty.
  \]

  $\Longrightarrow$ Không tồn tại $c_1 > 0$ thỏa mãn với \textbf{mọi} $n \geq n_0$. \textbf{Mâu thuẫn!}

  Vậy giả sử sai $\Rightarrow \frac{n(n+1)}{2} \notin \Theta(n^3)$.

  \bigskip

  \item $\dfrac{n(n+1)}{2} \in \Omega(n)$ \hfill \textcolor{darkgreen}{\textbf{ĐÚNG}} \\[4pt]
  \textbf{Lập luận:} $\frac{n(n+1)}{2} = \frac{n^2+n}{2} \geq \frac{n}{2}$ với mọi $n \geq 0$.

  Chọn $c = \frac{1}{2}$, $n_0 = 0$: ta có $\frac{n(n+1)}{2} \geq \frac{1}{2} \cdot n$ với mọi $n \geq 0$.

  $\Longrightarrow$ Theo định nghĩa $\Omega$: $\frac{n(n+1)}{2} \in \Omega(n)$. \checkmark

\end{enumerate}

\subsection{Câu 2: Xác định lớp $\Theta(g(n))$ đơn giản nhất}

\begin{enumerate}[label=\textbf{\alph*.}]

  \item $(n^2 + 1)^{10}$ \\[4pt]
  \textbf{Bước 1 --- Rút gọn:} Khai triển bậc cao nhất: $(n^2+1)^{10}$, hạng dẫn đầu là $n^{20}$.

  \textbf{Bước 2 --- Chứng minh $\Theta(n^{20})$:}

  Theo định nghĩa của $\Theta$, ta cần chỉ ra $\exists\, c_1, c_2 \in \mathbb{R}^+,\; n_0 \in \mathbb{N}$ sao cho:
  \[
    c_1 \cdot n^{20} \leq (n^2+1)^{10} \leq c_2 \cdot n^{20} \quad \forall\, n \geq n_0.
  \]

  \textit{Cận trên:} Với mọi $n \geq 1$ ta có $1 \leq n^2$, suy ra:
  \[
    (n^2 + 1)^{10} \leq (n^2 + n^2)^{10} = (2n^2)^{10} = 1024\, n^{20}.
  \]

  \textit{Cận dưới:} Với mọi $n \geq 0$ ta có:
  \[
    (n^2 + 1)^{10} \geq (n^2)^{10} = n^{20}.
  \]

  Suy ra chọn $c_1 = 1$, $c_2 = 1024$, $n_0 = 1$. Theo định nghĩa của $\Theta$ ta có ĐPCM.

  \begin{resultbox}
  $(n^2+1)^{10} \in \Theta(n^{20})$
  \end{resultbox}

  \bigskip

  \item $\sqrt{10n^2 + 7n + 3}$ \\[4pt]
  \textbf{Bước 1 --- Rút gọn:} Hạng chi phối dưới dấu căn là $10n^2$.

  \textbf{Bước 2 --- Chứng minh $\Theta(n)$:}

  Theo định nghĩa của $\Theta$, ta cần chỉ ra $\exists\, c_1, c_2 \in \mathbb{R}^+,\; n_0 \in \mathbb{N}$ sao cho:
  \[
    c_1 \cdot n \leq \sqrt{10n^2 + 7n + 3} \leq c_2 \cdot n \quad \forall\, n \geq n_0.
  \]

  \textit{Cận trên:} Với mọi $n \geq 1$ ta có $7n \leq 7n^2$ và $3 \leq 3n^2$, suy ra:
  \[
    \sqrt{10n^2 + 7n + 3} \leq \sqrt{10n^2 + 7n^2 + 3n^2} = \sqrt{20}\, n.
  \]

  \textit{Cận dưới:} Với mọi $n \geq 0$ ta có $7n + 3 \geq 0$, suy ra:
  \[
    \sqrt{10n^2 + 7n + 3} \geq \sqrt{10n^2} = \sqrt{10}\, n.
  \]

  Suy ra chọn $c_1 = \sqrt{10}$, $c_2 = \sqrt{20}$, $n_0 = 1$. Theo định nghĩa của $\Theta$ ta có ĐPCM.

  \begin{resultbox}
  $\sqrt{10n^2 + 7n + 3} \in \Theta(n)$
  \end{resultbox}

  \bigskip

  \item $2n\lg(n+2)^2 + (n+2)^2\lg\frac{n}{2}$ \\[4pt]
  \textbf{Bước 1 --- Rút gọn từng hạng:}
  \begin{itemize}[noitemsep]
    \item $2n\lg(n+2)^2 = 4n\lg(n+2) \in \Theta(n\log n)$.
    \item $(n+2)^2\lg\frac{n}{2} = (n+2)^2(\lg n - 1) \in \Theta(n^2 \log n)$.
  \end{itemize}

  \textbf{Bước 2 --- Hạng chi phối:} $n^2\log n$ thống trị $n\log n$ (bậc đa thức cao hơn).

  \begin{resultbox}
  $2n\lg(n+2)^2 + (n+2)^2\lg\frac{n}{2} \in \Theta(n^2\lg n)$
  \end{resultbox}

  \bigskip

  \item $2^{n+1} + 3^{n-1}$ \\[4pt]
  \textbf{Bước 1 --- Rút gọn:} $2^{n+1} = 2 \cdot 2^n$, $3^{n-1} = \frac{1}{3} \cdot 3^n$.

  \textbf{Bước 2 --- So sánh:} Đưa về cơ số 2: $2^n = 2^n$, $3^n = 2^{n\lg 3} \approx 2^{1{,}585\,n}$.

  Phần mũ $1{,}585\,n > n$ nên $3^n$ tăng nhanh hơn $2^n$. Hạng $3^{n-1}$ thống trị.

  \textbf{Bước 3 --- Chứng minh $\Theta(3^n)$:}

  Theo định nghĩa của $\Theta$, ta cần chỉ ra $\exists\, c_1, c_2 \in \mathbb{R}^+,\; n_0 \in \mathbb{N}$ sao cho:
  \[
    c_1 \cdot 3^n \leq 2^{n+1} + 3^{n-1} \leq c_2 \cdot 3^n \quad \forall\, n \geq n_0.
  \]

  \textit{Cận dưới:} Với mọi $n \geq 1$ ta có:
  \[
    2^{n+1} + 3^{n-1} \geq 3^{n-1} = \frac{1}{3} \cdot 3^n.
  \]

  \textit{Cận trên:} Với mọi $n \geq 1$ ta có $2^{n+1} = 2 \cdot 2^n \leq 2 \cdot 3^n$, suy ra:
  \[
    2^{n+1} + 3^{n-1} \leq 2 \cdot 3^n + \frac{1}{3} \cdot 3^n = \frac{7}{3} \cdot 3^n.
  \]

  Suy ra chọn $c_1 = 1/3$, $c_2 = 7/3$, $n_0 = 1$. Theo định nghĩa của $\Theta$ ta có ĐPCM.

  \begin{resultbox}
  $2^{n+1} + 3^{n-1} \in \Theta(3^n)$
  \end{resultbox}

  \bigskip

  \item $\lfloor\log_2 n\rfloor$ \\[4pt]
  \textbf{Chứng minh $\Theta(\log n)$:}

  Theo định nghĩa của $\Theta$, ta cần chỉ ra $\exists\, c_1, c_2 \in \mathbb{R}^+,\; n_0 \in \mathbb{N}$ sao cho:
  \[
    c_1 \cdot \log_2 n \leq \lfloor\log_2 n\rfloor \leq c_2 \cdot \log_2 n \quad \forall\, n \geq n_0.
  \]

  Dùng tính chất floor: $x - 1 < \lfloor x \rfloor \leq x$.

  \textit{Cận trên:} Với mọi $n \geq 1$ ta có:
  \[
    \lfloor\log_2 n\rfloor \leq \log_2 n = 1 \cdot \log_2 n.
  \]

  \textit{Cận dưới:} Với mọi $n \geq 4$ ta có $\log_2 n \geq 2$, suy ra:
  \[
    \lfloor\log_2 n\rfloor > \log_2 n - 1 \geq \log_2 n - \frac{1}{2}\log_2 n = \frac{1}{2}\log_2 n.
  \]

  Suy ra chọn $c_1 = 1/2$, $c_2 = 1$, $n_0 = 4$. Theo định nghĩa của $\Theta$ ta có ĐPCM.

  \begin{resultbox}
  $\lfloor\log_2 n\rfloor \in \Theta(\log n)$
  \end{resultbox}

\end{enumerate}

\subsection{Câu 3: Chứng minh mọi đa thức bậc $k$ thuộc $\Theta(n^k)$}

\textbf{Đề:} Cho $p(n) = a_k n^k + a_{k-1}n^{k-1} + \cdots + a_0$ với $a_k > 0$. CMR $p(n) \in \Theta(n^k)$.

Theo định nghĩa của $\Theta$, ta cần chỉ ra $\exists\, c_1, c_2 \in \mathbb{R}^+,\; n_0 \in \mathbb{N}$ sao cho:
\[
  c_1 \cdot n^k \leq p(n) \leq c_2 \cdot n^k \quad \forall\, n \geq n_0.
\]

\textbf{Bước 1 --- Cận trên $O(n^k)$:}

Với mọi $n \geq 1$ ta có $n^i \leq n^k$ ($\forall\, 0 \leq i \leq k$), suy ra:
\begin{align*}
  |p(n)| &\leq |a_k|n^k + |a_{k-1}|n^{k-1} + \cdots + |a_0| \\
         &\leq (|a_k| + |a_{k-1}| + \cdots + |a_0|) \cdot n^k \quad \forall\, n \geq 1.
\end{align*}
Suy ra chọn $c_2 = \sum_{i=0}^{k} |a_i|$.

\textbf{Bước 2 --- Cận dưới $\Omega(n^k)$:}
\[
  p(n) = n^k \left(a_k + \frac{a_{k-1}}{n} + \cdots + \frac{a_0}{n^k}\right)
\]
Khi $n \to \infty$, biểu thức trong ngoặc $\to a_k > 0$. Tồn tại $n_0 \in \mathbb{N}$ sao cho với mọi $n \geq n_0$, biểu thức $\geq \frac{a_k}{2}$.

Suy ra: $p(n) \geq \frac{a_k}{2} \cdot n^k$ với mọi $n \geq n_0$. Chọn $c_1 = \frac{a_k}{2}$.

Từ Bước 1 và Bước 2, theo định nghĩa của $\Theta$ ta có ĐPCM:
$p(n) \in \Theta(n^k)$. \hfill $\blacksquare$

% =========================================
%                BÀI 3
% =========================================
\newpage
\section{Bài 3: Chứng minh/Phản chứng các mệnh đề}

\begin{enumerate}[label=\textbf{\alph*.}]

  \item \textbf{Nếu $t(n) \in O(g(n))$ thì $g(n) \in \Omega(t(n))$} \hfill \textcolor{darkgreen}{\textbf{ĐÚNG}} \\[4pt]

  Giả thiết: $t(n) \leq c \cdot g(n)$ $\forall n \geq n_0$ (với $c > 0$).

  Chia hai vế cho $c$: $g(n) \geq \frac{1}{c} \cdot t(n)$ $\forall n \geq n_0$.

  Đặt $c' = 1/c > 0$. Theo định nghĩa $\Omega$: $g(n) \in \Omega(t(n))$. \hfill $\blacksquare$

  \bigskip

  \item $\Theta(\alpha g(n)) = \Theta(g(n))$ với $\alpha > 0$ \hfill \textcolor{darkgreen}{\textbf{ĐÚNG}} \\[4pt]

  Lấy $f(n) \in \Theta(\alpha g(n))$: $c_1 \alpha g(n) \leq f(n) \leq c_2 \alpha g(n)$.

  Đặt $c_1' = c_1\alpha$, $c_2' = c_2\alpha$ (đều dương): $c_1' g(n) \leq f(n) \leq c_2' g(n)$.

  $\Longrightarrow f(n) \in \Theta(g(n))$. Chiều ngược tương tự. \hfill $\blacksquare$

  \bigskip

  \item $\Theta(g(n)) = O(g(n)) \cap \Omega(g(n))$ \hfill \textcolor{darkgreen}{\textbf{ĐÚNG}} \\[4pt]

  ($\subseteq$): $f \in \Theta(g)$ $\Rightarrow$ $c_1 g \leq f \leq c_2 g$ $\Rightarrow$ $f \in O(g)$ và $f \in \Omega(g)$.

  ($\supseteq$): $f \in O(g) \cap \Omega(g)$ $\Rightarrow$ $f \leq c_2 g$ ($n \geq n_1$) và $f \geq c_1 g$ ($n \geq n_2$).

  Chọn $n_0 = \max(n_1, n_2)$: $c_1 g \leq f \leq c_2 g$ $\forall n \geq n_0$ $\Rightarrow$ $f \in \Theta(g)$. \hfill $\blacksquare$

  \bigskip

  \item \textbf{$\forall$ $t(n), g(n) \geq 0$: $t(n) \in O(g(n))$ hoặc $t(n) \in \Omega(g(n))$} \hfill \textcolor{red}{\textbf{SAI}} \\[4pt]

  \textbf{Phản ví dụ:}
  \[
    t(n) = \begin{cases} 1 & n \text{ lẻ} \\ n^2 & n \text{ chẵn} \end{cases}, \quad
    g(n) = \begin{cases} n^2 & n \text{ lẻ} \\ 1 & n \text{ chẵn} \end{cases}
  \]

  Với $n$ chẵn: $t(n)/g(n) = n^2 \to \infty$ $\Rightarrow$ không có $c$ cho $O$.

  Với $n$ lẻ: $t(n)/g(n) = 1/n^2 \to 0$ $\Rightarrow$ không có $c$ cho $\Omega$.

  $\Longrightarrow$ $t(n) \notin O(g(n))$ VÀ $t(n) \notin \Omega(g(n))$. \hfill $\blacksquare$

\end{enumerate}

\subsection{Exercises từ CLRS}

\textbf{3.1-1:} CMR $\max(f(n), g(n)) = \Theta(f(n) + g(n))$ với $f, g \geq 0$.

\textbf{Bước 1 --- Cận trên:} $\max(f,g) \leq f + g$. Chọn $c_2 = 1$. \checkmark

\textbf{Bước 2 --- Cận dưới:} $f + g \leq 2\max(f,g)$ $\Rightarrow$ $\max(f,g) \geq \frac{1}{2}(f+g)$. Chọn $c_1 = 1/2$. \checkmark

\begin{resultbox}
$\max(f(n), g(n)) = \Theta(f(n) + g(n))$ \hfill $\blacksquare$
\end{resultbox}

\bigskip

\textbf{3.1-2:} CMR $(n+a)^b = \Theta(n^b)$ với $b > 0$.

Khi $n \geq 2|a|$: $n/2 \leq n+a \leq 2n$

$\Rightarrow (n/2)^b \leq (n+a)^b \leq (2n)^b$, tức $\frac{1}{2^b} n^b \leq (n+a)^b \leq 2^b n^b$.

\begin{resultbox}
$(n+a)^b = \Theta(n^b)$ với $c_1 = 1/2^b$, $c_2 = 2^b$, $n_0 = 2|a|$. \hfill $\blacksquare$
\end{resultbox}

\bigskip

\textbf{3.1-4:} $2^{n+1} = O(2^n)$? \quad $2^{2n} = O(2^n)$?

$2^{n+1} = 2 \cdot 2^n \leq c \cdot 2^n$ với $c = 2$. $\Longrightarrow$ \textcolor{darkgreen}{\textbf{ĐÚNG}}. \checkmark

$2^{2n} = 4^n$. Giả sử $4^n \leq c \cdot 2^n$ $\Rightarrow$ $2^n \leq c$ $\Rightarrow$ mâu thuẫn vì $2^n \to \infty$.

$\Longrightarrow$ \textcolor{red}{\textbf{SAI}}. \hfill $\blacksquare$

% =========================================
%                BÀI 4
% =========================================
\newpage
\section{Bài 4: Asymptotic Notation Properties {\small\textcolor{red}{(không bắt buộc)}}}

\begin{enumerate}[label=\textbf{\alph*.}]

  \item $f(n) = O(g(n)) \Rightarrow g(n) = O(f(n))$ \hfill \textcolor{red}{\textbf{SAI}} \\[4pt]
  Phản ví dụ: $f(n) = n$, $g(n) = n^2$. $n \in O(n^2)$ nhưng $n^2 \notin O(n)$.

  \bigskip

  \item $f(n) + g(n) = \Theta(\min(f(n), g(n)))$ \hfill \textcolor{red}{\textbf{SAI}} \\[4pt]
  Phản ví dụ: $f(n) = n$, $g(n) = n^2$. $f+g = \Theta(n^2)$ nhưng $\min(f,g) = n$.

  Đúng ra: $f+g = \Theta(\max(f,g))$.

  \bigskip

  \item $f(n) = O(g(n)) \Rightarrow \lg f(n) = O(\lg g(n))$ khi $f,g \geq 2$ \hfill \textcolor{darkgreen}{\textbf{ĐÚNG}} \\[4pt]
  $f \leq cg$ $\Rightarrow$ $\lg f \leq \lg c + \lg g \leq (1+\lg c)\lg g$ khi $\lg g \geq 1$. \hfill $\blacksquare$

  \bigskip

  \item $f(n) = O(g(n)) \Rightarrow 2^{f(n)} = O(2^{g(n)})$ \hfill \textcolor{red}{\textbf{SAI}} \\[4pt]
  Phản ví dụ: $f(n) = 2n$, $g(n) = n$. $2n = O(n)$ nhưng $4^n \neq O(2^n)$.

  \bigskip

  \item $f(n) = O((f(n))^2)$ \hfill \textcolor{red}{\textbf{SAI}} \\[4pt]
  Phản ví dụ: $f(n) = 1/n$. $f^2 = 1/n^2$. Cần $1/n \leq c/n^2$ $\Rightarrow$ $n \leq c$. Không thể.

  \bigskip

  \item $f(n) = O(g(n)) \Rightarrow g(n) = \Omega(f(n))$ \hfill \textcolor{darkgreen}{\textbf{ĐÚNG}} \\[4pt]
  Giống bài 3a. $f \leq cg$ $\Rightarrow$ $g \geq (1/c)f$. \hfill $\blacksquare$

  \bigskip

  \item $f(n) = \Theta(f(n/2))$ \hfill \textcolor{red}{\textbf{SAI}} \\[4pt]
  Phản ví dụ: $f(n) = 2^n$. $f(n)/f(n/2) = 2^{n/2} \to \infty$.

  \bigskip

  \item $f(n) + o(f(n)) = \Theta(f(n))$ \hfill \textcolor{darkgreen}{\textbf{ĐÚNG}} \\[4pt]
  Gọi $h = o(f)$, tức $h/f \to 0$.
  \begin{itemize}[noitemsep]
    \item Cận trên: $f + h \leq f + f = 2f$ khi $n$ đủ lớn.
    \item Cận dưới: $f + h \geq f - |h| \geq f/2$ khi $n$ đủ lớn.
  \end{itemize}
  $\Longrightarrow f + o(f) = \Theta(f)$. \hfill $\blacksquare$

\end{enumerate}

% =========================================
%                BÀI 5
% =========================================
\newpage
\section{Bài 5: Tính tổng}

\subsection{Phần 1: Compute the following sums}

\begin{enumerate}[label=\textbf{\alph*.}]

  \item $1 + 3 + 5 + 7 + \cdots + 999$ \\[4pt]
  Cấp số cộng: $a_1 = 1$, $d = 2$, $a_n = 999$.

  Số phần tử: $999 = 1 + (n-1) \cdot 2$ $\Rightarrow$ $n = 500$.

  $S = \dfrac{n(a_1 + a_n)}{2} = \dfrac{500 \times 1000}{2}$.

  \begin{resultbox}
  $S = 250\,000$
  \end{resultbox}

  \bigskip

  \item $2 + 4 + 8 + 16 + \cdots + 1024$ \\[4pt]
  Cấp số nhân: $a = 2$, $r = 2$, $ar^{n-1} = 1024$ $\Rightarrow$ $n = 10$.

  $S = \dfrac{a(r^n - 1)}{r - 1} = \dfrac{2(2^{10} - 1)}{1}$.

  \begin{resultbox}
  $S = 2046$
  \end{resultbox}

  \bigskip

  \item $\displaystyle\sum_{i=3}^{n+1} 1$ \\[4pt]
  Đếm: $i$ chạy từ $3$ đến $n+1$, có $(n+1) - 3 + 1 = n - 1$ phần tử.

  \begin{resultbox}
  $\displaystyle\sum_{i=3}^{n+1} 1 = n - 1$
  \end{resultbox}

  \bigskip

  \item $\displaystyle\sum_{i=3}^{n+1} i$ \\[4pt]
  Tách: $\displaystyle\sum_{i=3}^{n+1} i = \sum_{i=1}^{n+1} i - (1+2) = \frac{(n+1)(n+2)}{2} - 3$.

  \begin{resultbox}
  $\displaystyle\sum_{i=3}^{n+1} i = \frac{(n+4)(n-1)}{2}$
  \end{resultbox}

  \bigskip

  \item $\displaystyle\sum_{i=0}^{n-1} i(i+1)$ \\[4pt]
  Khai triển: $i(i+1) = i^2 + i$.
  \[
    \sum_{i=0}^{n-1}(i^2 + i) = \frac{(n-1)n(2n-1)}{6} + \frac{(n-1)n}{2}
    = \frac{(n-1)n(2n+2)}{6}
  \]

  \begin{resultbox}
  $\displaystyle\sum_{i=0}^{n-1} i(i+1) = \frac{n(n-1)(n+1)}{3}$
  \end{resultbox}

  \bigskip

  \item $\displaystyle\sum_{j=1}^{n} 3^{j+1}$ \\[4pt]
  $= 3\displaystyle\sum_{j=1}^{n} 3^j = 3 \cdot \frac{3(3^n - 1)}{2}$.

  \begin{resultbox}
  $\displaystyle\sum_{j=1}^{n} 3^{j+1} = \frac{9(3^n - 1)}{2}$
  \end{resultbox}

  \bigskip

  \item $\displaystyle\sum_{i=1}^{n}\sum_{j=1}^{n} ij$ \\[4pt]
  Tách: $= \left(\displaystyle\sum_{i=1}^{n} i\right) \cdot \left(\sum_{j=1}^{n} j\right) = \left(\frac{n(n+1)}{2}\right)^2$.

  \begin{resultbox}
  $\displaystyle\sum_{i=1}^{n}\sum_{j=1}^{n} ij = \frac{n^2(n+1)^2}{4}$
  \end{resultbox}

  \bigskip

  \item $\displaystyle\sum_{i=1}^{n} \frac{1}{i(i+1)}$ \\[4pt]
  Phân tích: $\frac{1}{i(i+1)} = \frac{1}{i} - \frac{1}{i+1}$ (tổng telescoping).
  \[
    = \left(1 - \frac{1}{2}\right) + \left(\frac{1}{2} - \frac{1}{3}\right) + \cdots + \left(\frac{1}{n} - \frac{1}{n+1}\right) = 1 - \frac{1}{n+1}
  \]

  \begin{resultbox}
  $\displaystyle\sum_{i=1}^{n} \frac{1}{i(i+1)} = \frac{n}{n+1}$
  \end{resultbox}

\end{enumerate}

\subsection{Phần 2: Tìm bậc tăng trưởng $\Theta(g(n))$}

\begin{enumerate}[label=\textbf{\alph*.}]

  \item $\displaystyle\sum_{i=0}^{n-1}(i^2+1)^2$ \\[4pt]
  Khai triển: $(i^2+1)^2 = i^4 + 2i^2 + 1$. Hạng chi phối: $\sum i^4 = \Theta(n^5)$.

  \begin{resultbox}
  $\displaystyle\sum_{i=0}^{n-1}(i^2+1)^2 \in \Theta(n^5)$
  \end{resultbox}

  \bigskip

  \item $\displaystyle\sum_{i=2}^{n-1}\lg i^2$ \\[4pt]
  $\lg i^2 = 2\lg i$. $\sum 2\lg i = \Theta(\lg(n!)) = \Theta(n\log n)$ (Stirling).

  \begin{resultbox}
  $\displaystyle\sum_{i=2}^{n-1}\lg i^2 \in \Theta(n\log n)$
  \end{resultbox}

  \bigskip

  \item $\displaystyle\sum_{i=1}^{n}(i+1)2^{i-1}$ \\[4pt]
  Dạng $\sum i \cdot 2^i$ có closed-form $= \Theta(n \cdot 2^n)$.

  \begin{resultbox}
  $\displaystyle\sum_{i=1}^{n}(i+1)2^{i-1} \in \Theta(n \cdot 2^n)$
  \end{resultbox}

  \bigskip

  \item $\displaystyle\sum_{i=0}^{n-1}\sum_{j=0}^{i-1}(i+j)$ \\[4pt]
  Tổng trong: $\sum_{j=0}^{i-1}(i+j) = i^2 + \frac{i(i-1)}{2} = \frac{3i^2 - i}{2}$.

  Tổng ngoài: $\sum \frac{3i^2-i}{2}$. Hạng chi phối: $\frac{3}{2}\sum i^2 = \Theta(n^3)$.

  \begin{resultbox}
  $\displaystyle\sum_{i=0}^{n-1}\sum_{j=0}^{i-1}(i+j) \in \Theta(n^3)$
  \end{resultbox}

\end{enumerate}

% =========================================
%                BÀI 6
% =========================================
\newpage
\section{Bài 6: Comparison of Running Times {\small\textcolor{red}{(không bắt buộc)}}}

\textbf{Phương pháp:} Giải $f(n) \leq t$ (đổi $t$ ra micro giây).

\begin{notebox}
\begin{center}
\renewcommand{\arraystretch}{1.2}
\begin{tabular}{|l|r|}
\hline
\textbf{Thời gian} & \textbf{Micro giây} \\
\hline
1 second & $10^6$ \\
1 minute & $6 \times 10^7$ \\
1 hour & $3.6 \times 10^9$ \\
1 day & $8.64 \times 10^{10}$ \\
1 month & $\approx 2.592 \times 10^{12}$ \\
1 year & $\approx 3.154 \times 10^{13}$ \\
1 century & $\approx 3.154 \times 10^{15}$ \\
\hline
\end{tabular}
\end{center}
\end{notebox}

\textbf{Ví dụ với $t = 1$ second $= 10^6$ $\mu$s:}

\begin{center}
\renewcommand{\arraystretch}{1.4}
\begin{tabular}{|c|l|c|}
\hline
$f(n)$ & \textbf{Cách giải} & $n$ max \\
\hline
$\lg n$ & $\lg n \leq 10^6 \Rightarrow n \leq 2^{10^6}$ & $2^{10^6}$ \\
\hline
$\sqrt{n}$ & $\sqrt{n} \leq 10^6 \Rightarrow n \leq 10^{12}$ & $10^{12}$ \\
\hline
$n$ & $n \leq 10^6$ & $10^6$ \\
\hline
$n\lg n$ & Xấp xỉ: $n\lg n \approx 10^6$ & $\approx 6.27 \times 10^4$ \\
\hline
$n^2$ & $n^2 \leq 10^6 \Rightarrow n \leq 1000$ & $1000$ \\
\hline
$n^3$ & $n^3 \leq 10^6 \Rightarrow n \leq 100$ & $100$ \\
\hline
$2^n$ & $2^n \leq 10^6 \Rightarrow n \leq 19.93$ & $19$ \\
\hline
$n!$ & Thử: $9! = 362880 < 10^6$, $10! > 10^6$ & $9$ \\
\hline
\end{tabular}
\end{center}

Áp dụng tương tự cho các cột khác bằng cách thay $10^6$ bằng giá trị tương ứng.

% =========================================
%                BÀI 7
% =========================================
\newpage
\section{Bài 7: Relative Asymptotic Growths {\small\textcolor{red}{(không bắt buộc)}}}

\begin{notebox}
$A \in O(B)$: $A$ tăng $\leq B$. \quad $A \in o(B)$: $A$ tăng chặt $< B$.

$A \in \Omega(B)$: $A$ tăng $\geq B$. \quad $A \in \omega(B)$: $A$ tăng chặt $> B$. \quad $A \in \Theta(B)$: cùng bậc.
\end{notebox}

\begin{enumerate}[label=\textbf{\alph*.}]

  \item $A = \lg^k n$, $B = n^\epsilon$ ($k \geq 1$, $\epsilon > 0$) \\[4pt]
  Mọi $\log^k n$ tăng chậm hơn $n^\epsilon$: $\lim \frac{\lg^k n}{n^\epsilon} = 0$.

  \begin{center}
  \begin{tabular}{|c|c|c|c|c|}
  \hline $O$ & $o$ & $\Omega$ & $\omega$ & $\Theta$ \\ \hline
  yes & yes & no & no & no \\ \hline
  \end{tabular}
  \end{center}

  \item $A = n^k$, $B = c^n$ ($k \geq 1$, $c > 1$) \\[4pt]
  Đa thức luôn tăng chậm hơn hàm mũ: $\lim \frac{n^k}{c^n} = 0$.

  \begin{center}
  \begin{tabular}{|c|c|c|c|c|}
  \hline $O$ & $o$ & $\Omega$ & $\omega$ & $\Theta$ \\ \hline
  yes & yes & no & no & no \\ \hline
  \end{tabular}
  \end{center}

  \item $A = \sqrt{n}$, $B = n^{\sin n}$ \\[4pt]
  $\sin n$ dao động trong $[-1,1]$. Khi $\sin n \approx 0$: $B \approx 1 \ll A$. Khi $\sin n = 1$: $B = n \gg A$.

  Không có quan hệ tiệm cận ổn định.

  \begin{center}
  \begin{tabular}{|c|c|c|c|c|}
  \hline $O$ & $o$ & $\Omega$ & $\omega$ & $\Theta$ \\ \hline
  no & no & no & no & no \\ \hline
  \end{tabular}
  \end{center}

  \item $A = 2^n$, $B = 2^{n/2}$ \\[4pt]
  $A/B = 2^{n/2} \to \infty$. $A$ tăng nhanh hơn $B$.

  \begin{center}
  \begin{tabular}{|c|c|c|c|c|}
  \hline $O$ & $o$ & $\Omega$ & $\omega$ & $\Theta$ \\ \hline
  no & no & yes & yes & no \\ \hline
  \end{tabular}
  \end{center}

  \item $A = n^{\lg c}$, $B = c^{\lg n}$ \\[4pt]
  Lấy $\lg$ hai vế: $\lg A = \lg c \cdot \lg n = \lg B$. Vậy $A = B$.

  \begin{center}
  \begin{tabular}{|c|c|c|c|c|}
  \hline $O$ & $o$ & $\Omega$ & $\omega$ & $\Theta$ \\ \hline
  yes & no & yes & no & yes \\ \hline
  \end{tabular}
  \end{center}

  \item $A = \lg(n!)$, $B = \lg(n^n)$ \\[4pt]
  $\lg(n^n) = n\lg n$. $\lg(n!) = \Theta(n\lg n)$ (Stirling). $\Longrightarrow A = \Theta(B)$.

  \begin{center}
  \begin{tabular}{|c|c|c|c|c|}
  \hline $O$ & $o$ & $\Omega$ & $\omega$ & $\Theta$ \\ \hline
  yes & no & yes & no & yes \\ \hline
  \end{tabular}
  \end{center}

\end{enumerate}

% =========================================
%                BÀI 8
% =========================================
\newpage
\section{Bài 8: Đoán lớp hiệu quả từ dữ liệu thực nghiệm {\small\textcolor{red}{(không bắt buộc)}}}

\textbf{Dữ liệu:}
\begin{center}
\renewcommand{\arraystretch}{1.1}
\begin{tabular}{|c|c|c|c|c|c|c|c|c|c|c|}
\hline
\textbf{size} & 1000 & 2000 & 3000 & 4000 & 5000 & 6000 & 7000 & 8000 & 9000 & 10000 \\
\hline
\textbf{count} & 11966 & 24303 & 39992 & 53010 & 67272 & 78692 & 91274 & 113063 & 129799 & 140538 \\
\hline
\end{tabular}
\end{center}

\textbf{Bước 1 --- Thử $\Theta(n\log n)$:} Tính $\frac{\text{count}}{n\lg n}$:

\begin{center}
\begin{tabular}{|c|c|c|c|}
\hline
$n$ & $n\lg n$ & count & $\frac{\text{count}}{n\lg n}$ \\
\hline
1000 & $\approx 9966$ & 11966 & $\approx 1.20$ \\
2000 & $\approx 21932$ & 24303 & $\approx 1.11$ \\
5000 & $\approx 61439$ & 67272 & $\approx 1.09$ \\
10000 & $\approx 132877$ & 140538 & $\approx 1.06$ \\
\hline
\end{tabular}
\end{center}

Tỷ số gần hằng $\approx 1.1$ và ổn định khi $n$ tăng. \checkmark

\textbf{Bước 2 --- Thử $\Theta(n^2)$:} $\frac{11966}{10^6} = 0.012$, $\frac{140538}{10^8} = 0.0014$. Giảm mạnh $\Rightarrow$ loại.

\begin{resultbox}
Lớp hiệu quả: $\Theta(n\log n)$.
\end{resultbox}

% =========================================
%                BÀI 9
% =========================================
\newpage
\section{Bài 9: Giải hệ thức truy hồi {\small\textcolor{red}{(không bắt buộc)}}}

\subsection{9.1: Backward substitution}

\textbf{a)} $T(n) = 2T(n/2) + \log n$, $T(1) = 1$

Giả sử $n = 2^k$. Khai triển ngược:
\begin{align*}
  T(n) &= 2T(n/2) + \log n \\
       &= 4T(n/4) + 2\log(n/2) + \log n \\
       &= 2^k T(1) + \sum_{i=0}^{k-1} 2^i(k-i) = n + 2^k\sum_{j=1}^{k}\frac{j}{2^j}
\end{align*}

Dùng $\sum_{j=1}^{\infty} j/2^j = 2$: tổng $\leq 2n$.

\begin{resultbox}
$T(n) = \Theta(n)$
\end{resultbox}

\bigskip

\textbf{b)} $T(n) = 4T(n/3) + n$, $T(1) = 1$

Giả sử $n = 3^k$. Khai triển:
\[
  T(n) = 4^k + n\sum_{i=0}^{k-1}(4/3)^i = 4^k + 3n[(4/3)^k - 1]
\]

Vì $4^k = 4^{\log_3 n} = n^{\log_3 4} \approx n^{1.26}$. Hạng $4^k$ thống trị.

\begin{resultbox}
$T(n) = \Theta(n^{\log_3 4}) \approx \Theta(n^{1.26})$
\end{resultbox}

\subsection{9.2: Phương trình đặc trưng}

\textbf{a)} $T(n) = 7T(n-1) - 12T(n-2)$, $T(0) = 1$, $T(1) = 2$

PT đặc trưng: $x^2 - 7x + 12 = 0$ $\Rightarrow$ $x_1 = 4$, $x_2 = 3$.

Nghiệm tổng quát: $T(n) = A \cdot 4^n + B \cdot 3^n$.

Điều kiện đầu: $A + B = 1$, $4A + 3B = 2$ $\Rightarrow$ $A = -1$, $B = 2$.

\begin{resultbox}
$T(n) = 2 \cdot 3^n - 4^n$

Kiểm tra: $T(0) = 2 - 1 = 1$ \checkmark, $T(1) = 6 - 4 = 2$ \checkmark
\end{resultbox}

\bigskip

\textbf{b)} $T(n+1) = T(n) + 3n$, $T(0) = 7$

Khai triển: $T(n) = T(0) + 3\sum_{i=0}^{n-1}i = 7 + \frac{3n(n-1)}{2}$.

\begin{resultbox}
$T(n) = \dfrac{3n^2 - 3n + 14}{2}$

Kiểm tra: $T(0) = 7$ \checkmark, $T(1) = 7$ \checkmark, $T(2) = 10 = T(1) + 3$ \checkmark
\end{resultbox}

\bigskip

\textbf{c)} $T(n) = 4T(n-1) - 5T(n-2) + 2T(n-3)$, $T(0)=0$, $T(1)=1$, $T(2)=2$

PT đặc trưng: $x^3 - 4x^2 + 5x - 2 = 0$. Thử $x=1$: $1-4+5-2=0$ \checkmark.

Chia Horner: $(x-1)^2(x-2) = 0$. Nghiệm kép $x=1$ (bội 2), $x=2$.

$T(n) = A + Bn + C \cdot 2^n$. Hệ: $A+C = 0$, $A+B+2C = 1$, $A+2B+4C = 2$.

Giải: $C = 0$, $B = 1$, $A = 0$.

\begin{resultbox}
$T(n) = n$

Kiểm tra: $T(3) = 4(2)-5(1)+2(0) = 3$ \checkmark
\end{resultbox}

\bigskip

\textbf{d)} $T(n) = T(n-1) + T(n-2)$, $T(0) = 1$, $T(1) = 1$ (dạng Fibonacci)

PT đặc trưng: $x^2 - x - 1 = 0$ $\Rightarrow$ $x = \frac{1 \pm \sqrt{5}}{2}$.

Đặt $\phi = \frac{1+\sqrt{5}}{2} \approx 1.618$, $\hat\phi = \frac{1-\sqrt{5}}{2} \approx -0.618$.

\begin{resultbox}
$T(n) = \dfrac{\phi^{n+1} - \hat\phi^{n+1}}{\sqrt{5}} = \Theta(\phi^n) \approx \Theta(1.618^n)$
\end{resultbox}

\subsection{9.3: Đoán nghiệm (Substitution method)}

\textbf{a)} $T(n) = 3T(n/2) + n^2$, $T(1) = 1$

Dùng Master Theorem: $a = 3$, $b = 2$, $f(n) = n^2$.

$n^{\log_2 3} \approx n^{1.585}$. Mà $f(n) = n^2 = \Omega(n^{1.585+\epsilon})$ $\Rightarrow$ Case 3.

Regularity: $3(n/2)^2 = 3n^2/4 \leq cn^2$ với $c = 3/4 < 1$ \checkmark.

\begin{resultbox}
$T(n) = \Theta(n^2)$
\end{resultbox}

\bigskip

\textbf{b)} $T(n) = T(n/2) + T(n/4) + n$, $T(n) = 1$ ($n \leq 5$). Gợi ý: $T(n) = O(n)$.

Đoán $T(n) \leq cn$. Thay vào:
\[
  T(n) \leq cn/2 + cn/4 + n = n(3c/4 + 1)
\]

Cần $3c/4 + 1 \leq c$ $\Rightarrow$ $c \geq 4$. Base: $T(n) = 1 \leq 4n$ ($n \geq 1$) \checkmark.

\begin{resultbox}
$T(n) = O(n)$ với $c = 4$.
\end{resultbox}

% =========================================
%                TỔNG KẾT
% =========================================
\newpage
\section*{Tổng kết phương pháp}

\begin{notebox}
\begin{enumerate}
  \item \textbf{Sắp xếp hàm:} Phân lớp $\to$ So sánh bằng tính chất/cùng cơ số $\to$ Sắp xếp
  \item \textbf{True/False:} Dùng định nghĩa formal, chọn $c$ và $n_0$ cụ thể
  \item \textbf{Chứng minh:} Tìm $c_1, c_2, n_0$; phản chứng thì cho counterexample
  \item \textbf{Tính tổng:} CSC, CSN, telescoping, tách tổng
  \item \textbf{Truy hồi:} Backward substitution, PT đặc trưng, Master Theorem, Đoán nghiệm
\end{enumerate}
\end{notebox}

\end{document}
